% ===== 原文 pp.4--5 =====
% tr-01 片 a：§2 标题、§2.1 标题、§2.1.1 作用量、§2.1.2 在圆柱面上、§2.1.3 前半
% 编号公式 tags: 2.1--2.9
% 记号说明：原文的度量写成无衬线体的 g（PDF 内嵌字体为 CMSS），用来与群元素 $g$ 区分，
% 本译文照排为 \mathsf{g}。

% [原文 p.4]
\section{两个二维理论}

\subsection{二维杨--米尔斯理论}

我们首先研究二维杨--米尔斯理论（two-dimensional Yang-Mills theories）。我们先处理标准的、
未作形变的规范理论，并在最后说明如何对它作 $q$-形变。未形变理论的每个细节都可以在
出色的综述 [10] 中找到。

\subsubsection{作用量}

规范群为 $G$ 的二维杨--米尔斯理论具有拉格朗日量
\begin{equation}
  S \propto \frac{1}{e^2}\int d^2x\,\sqrt{\det\mathsf{g}}\,\tr F_{\mu\nu}F^{\mu\nu}.
  \tag{2.1}
\end{equation}
这里我们考虑的是带背景度规 $\mathsf{g}$ 的弯曲流形上的理论，采用欧几里得号差。耦合常数 $e$ 可以
通过重新标度 $\mathsf{g}$ 去掉。

回想一下，在二维情形，场强 $F_{\mu\nu}$ 唯一非零的分量是 $F_{01}$。于是动能项可以写成
\begin{equation}
  \tr F_{\mu\nu}F^{\mu\nu}
  \propto (\mathsf{g}^{00}\mathsf{g}^{11}-\mathsf{g}^{01}\mathsf{g}^{10})\,\tr(F_{01})^2
  = (\det\mathsf{g})^{-1}\tr(F_{01})^2 .
  \tag{2.2}
\end{equation}
% [Tutor Note: p.4, 式 (2.2)：原文把中间步骤压成了一步，补出来是
%   $F^{01}=\mathsf{g}^{0\mu}\mathsf{g}^{1\nu}F_{\mu\nu}
%   =\mathsf{g}^{00}\mathsf{g}^{11}F_{01}+\mathsf{g}^{01}\mathsf{g}^{10}F_{10}
%   =(\mathsf{g}^{00}\mathsf{g}^{11}-\mathsf{g}^{01}\mathsf{g}^{10})F_{01}$（用了 $F_{10}=-F_{01}$），
%   于是 $\tr F_{\mu\nu}F^{\mu\nu}=2\tr(F_{01}F^{01})
%   =2(\mathsf{g}^{00}\mathsf{g}^{11}-\mathsf{g}^{01}\mathsf{g}^{10})\tr(F_{01})^2$，
%   其中因子 $2$ 被 $\propto$ 吸收；最后一步用
%   $\mathsf{g}^{00}\mathsf{g}^{11}-\mathsf{g}^{01}\mathsf{g}^{10}=\det(\mathsf{g}^{-1})=(\det\mathsf{g})^{-1}$。
%   易错点：原文的 $\tr(F_{01})^2$ 指的是 $\tr\bigl((F_{01})^2\bigr)$，即先把同一分量的场强作矩阵
%   平方再取迹；它不是 $(\tr F_{01})^2$。原文省掉了这层括号，而这两种读法只有前者正确。]
这意味着在作用量 (2.1) 中并不出现度规 $\mathsf{g}$ 的各个分量：唯一出现的组合是
$\det\mathsf{g}$。换个说法，二维杨--米尔斯理论可以建立在并不真正带有「能测量距离」的度规的
二维曲面上；我们需要的全部只是一个体积形式 $dx^0dx^1\sqrt{\det\mathsf{g}}$，它使我们能够
测量面积。于是二维曲面的唯一不变量就是它的亏格与曲面的总面积，配分函数只能依赖于这两者。

\subsubsection{在圆柱面上}

现在我们来分析圆柱面 $x^1\sim x^1+L$ 上的理论，取 $x^0$ 为时间方向，见图 1。我们取时间规范
$A_0=0$。在等时切片 $x^0=0$ 上，规范不变的数据是规范场绕该圆的和乐（holonomy）：
\begin{equation}
  U := P\exp\int_0^L A_1dx^1 \in G,
  \tag{2.3}
\end{equation}

% [原文 p.5]
\begin{center}
\begin{minipage}{\textwidth}
\centering
\includegraphics[width=0.35\textwidth]{fig1.png}\par\vspace{2pt}
\figtitle{1}{一个周长为 $L$ 的圆柱面}
\end{minipage}
\end{center}
其中 $U$ 在伴随 $G$-作用 $U\mapsto gUg^{-1}$ 下视为等价。

于是系统的波函数是定义在 $G$ 上的函数 $\psi(U)$，它满足不变性
\begin{equation}
  \psi(U)=\psi(gUg^{-1}).
  \tag{2.4}
\end{equation}
在群论中这是一个标准事实：这类函数由不可约表示 $R$ 中迹的线性组合给出
\begin{equation}
  \chi_R(U)=\tr_R U
  \tag{2.5}
\end{equation}
注意它们在 $G$ 上的自然测度下正交归一：
\begin{equation}
  \int_G\chi_R(U)\chi_{R'}(U)^*\,dU=\delta_{RR'}.
  \tag{2.6}
\end{equation}

把 $A_1$ 分解为 $A_1=A_1^aT_a$，其中 $a=1,\dots,\dim G$，由 (2.1) 得到的哈密顿量为
\begin{equation}
  H \propto \int_0^L \frac{\delta}{\delta A_1^a(x)}\frac{\delta}{\delta A_1^a(x)}\,dx^1 .
  \tag{2.7}
\end{equation}
% [Tutor Note: p.5, 式 (2.7)：由 (2.1) 到 (2.7) 之间省略了正则量子化。在 $A_0=0$ 下
%   $L\propto\frac{1}{e^2}\int dx^1\,\tr(\dot A_1)^2$，共轭动量 $\Pi^a\propto\dot A_1^a/e^2$，
%   于是 $H\propto e^2\int dx^1\,(\Pi^a)^2$；把 $\Pi^a$ 换成 $-i\delta/\delta A_1^a(x)$（$\hbar=1$），
%   再把整体常数并入 $\propto$，就是 (2.7)。
%   易错点：两个泛函导数都作用在同一个 $A_1^a(x)$ 上（同一坐标 $x$，不是 $x'$），
%   这种「同点乘积」在排版上极容易看漏；写成不同的 $x,x'$ 再取极限，结论相同。]
作用在 $\chi_R(U)=\tr_R P\exp\int_0^L A_1dx^1$ 上，我们得到
\begin{equation}
  H\chi_R(U)\propto \tr_R\int_0^L T^aT^a dx^1\,P\exp\Bigl(\int_0^L A_1dx^1\Bigr)
  \propto Lc_2(R)\chi_R(U)
  \tag{2.8}
\end{equation}
% [Tutor Note: p.5, 式 (2.8) 的两个 $\propto$：泛函导数作用在 $\chi_R(U)=\tr_R P\exp\int_0^L A_1dx^1$
%   上时，每次作用都在路径有序指数中插入一个 $T^a$（插入点即 $x$）；两次作用给出 $T^aT^a$，
%   而在不可约表示 $R$ 中 $T^aT^a=c_2(R)\mathbf{1}$，于是得到 $c_2(R)\chi_R(U)$。
%   中间的因子 $L$ 来自 $\int_0^L dx^1$：插入点扫过整个圆，给出圆周长 $L$。
%   易错点：这个 $L$ 与 $\dim R$ 无关，不要以为它来自 $\tr_R\mathbf{1}$；最终 (2.9) 里也确无 $\dim R$。]
其中 $c_2(R)$ 是二次 Casimir 算子在不可约表示 $R$ 中的取值。我们通过要求
$H\chi_R=Lc_2(R)\chi_R$ 来确定各比例常数。

现在可以计算环面 $x^1\sim x^1+L$、$x^0\sim x^1+T$ 上的配分函数 $Z$：
\begin{equation}
  Z=\tr e^{-TH}=\sum_R e^{-TLc_2(R)} .
  \tag{2.9}
\end{equation}
% [Tutor Note: p.5, 式 (2.9) 前一句：原文把环面的两个周期条件写成 $x^1\sim x^1+L$ 与 $x^0\sim x^1+T$，
%   第二个条件里的 $x^1$ 应为 $x^0$（已按 700 dpi 裁图复核印刷原文，确系如此）。
%   正确写法是 $x^0\sim x^0+T$；照字面理解，$x^0\sim x^1+T$ 并不是任何环面的周期条件。
%   正文照原文排印；读者按 (2.9) 复算时请用 $x^0\sim x^0+T$（此时总面积确为 $TL$）。]
注意最终结果只依赖于环面的总面积 $TL$，这符合预期。

\subsubsection{在一般曲面上}

接下来，我们研究一般二维曲面上的理论。先考虑面积为 $A$ 的圆盘，见图 2。我们在边界圆上指定
和乐 $U$；然后可以在这一条件下计算配分函数。可以把它记作 $Z_A(U)$。这也可以看作在边界
$S^1$ 上定义了一个波函数，并且可以记作 $\psi_A(U)=Z_A(U)$。
